\begin{proposition}[ARI \ref{ari:C.6.R3}; Proved: the consequences of Gross--Zagier and Kolyvagin]\label{ari:C.6.R3}
For each of the nine rank-one curves there is a positive rational
number $q_i$ such that
\begin{equation}\label{ari:eq:Cr-rational-factor}
 \frac{L'(E_i,1)}{\Omega_i}=q_i\widehat h(P_i).
\end{equation}
Their Tate--Shafarevich groups over $\Q$ are finite. Consequently
the corresponding $D_{\rm part}$ is an explicit logarithm minus
$(6/(\pi q_i))L'(E_i,1)/\Omega_i$. For the two rank-two curves, $L'(E_i,1)=0$ while
$\widehat h(P_i)>0$.
\end{proposition}
\Uses{the hypotheses, notation and definitions stated here; \citep{CrGrossZagier1986}; \citet[Theorem 1.3]{CrGross1991}; \citet[Theorem 1.1]{CrJetchev2008}}

\begin{proof}
\begin{enumerate}[label=\arabic*.,leftmargin=*]
\item Gross--Zagier's theorem applies because each curve is modular and
$L(E_i,1)=0$. CM modularity follows from its Hecke-character
theta series. Their Theorem I.7.3, in the opened original
preprint \citep{CrGrossZagier1986}, gives a rational point $R$ and
a nonzero rational factor relating $L'(E,1)$ to its canonical
height times the real period. The differentials in the two conventions are rational multiples:
write $\omega'=a\omega_E$ with $a\in\Q^*$.  A primitive real
period for $\omega'$ is $|a|\Omega_i/n_{\R}$, where
$n_{\R}\in\{1,2\}$ is the number of real components.
Thus the change to the full real period multiplies the rational
factor by $|a|/n_{\R}$.
As $L'(E_i,1)>0$, $R$ has infinite order. Primitivity gives
$R=mP_i+T$, $m\ne0$, $T$ torsion; its height is
$m^2\widehat h(P_i)$. This proves
\eqref{ari:eq:Cr-rational-factor}, including positivity.

\item The rank-one Gross--Zagier--Kolyvagin implication asserts rank one
and finite Tate--Shafarevich group when the central zero is simple.
One can use Kolyvagin's Heegner-point theorem as stated in
\citet[Theorem 1.3]{CrGross1991}, together with the nonzero
Gross--Zagier Heegner height; a modern primary research statement
of the rank conclusion is \citet[Theorem 1.1]{CrJetchev2008}.
The hypotheses are satisfied separately by the positive derivative
and exact sign in each of the nine rows.  Explicitly, the auxiliary
imaginary quadratic field $K$ in the proof of Gross--Zagier I.7.3
has split Heegner primes and $L(E_i^D,1)\ne0$; hence
$L'(E_i/K,1)=L'(E_i,1)L(E_i^D,1)\ne0$.  Its Heegner point has
positive height, so Gross's Theorem~1.3 gives finite
$\operatorname{Sha}(E_i/K)$.  The restriction map
\[
\operatorname{Sha}(E_i/\Q)\longrightarrow
\operatorname{Sha}(E_i/K)
\]
has kernel contained in
$H^1(\operatorname{Gal}(K/\Q),E_i(K))$ by restriction--inflation.
This last group is finite: the Mordell--Weil group is finitely
generated, and its degree-one cohomology for this order-two group
is finitely generated and killed by two.  The finite kernel and
finite image prove the claimed finiteness over $\Q$.
For the two other rows the established equations are
$L'(E_i,1)=0$ and $\widehat h_O(P_i)>0$; therefore an equation
$L'(E_i,1)/\Omega_i=q\widehat h_O(P_i)$ with $q>0$ would contradict
these exact values.  Their rank-two regulator, for any free basis
$P,Q$, is the determinant
$\widehat h_O(P)\widehat h_O(Q)-\langle P,Q\rangle^2$.
\end{enumerate}
\end{proof}